\capituloUDA{Metodología}{Proceso de modelación, simulación y evaluación}{%
  \itemCapituloUDA{-2.15}{Diseño y enfoque de la investigación}
  \itemCapituloUDA{-3.18}{Caso de estudio y unidad de análisis}
  \itemCapituloUDA{-4.21}{Fases metodológicas}
  \itemCapituloUDA{-5.24}{Fuentes de información y datos de entrada}
  \itemCapituloUDA{-6.27}{Herramientas de simulación y análisis}
  \itemCapituloUDA{-7.30}{Escenarios de simulación e indicadores}
  \itemCapituloUDA{-8.33}{Procedimiento y validación del modelo}
}



\vspace*{1.15cm}
\section{Diseño y enfoque de la investigación}

\begin{textoDosTerciosGrafico}{\graficoDisenoEnfoqueMetodologia}{Diseño y enfoque metodológico de la investigación}{Diagrama original de la investigación.}
La metodología define el procedimiento para responder a la pregunta de investigación: ¿qué impacto tienen las estrategias pasivas y el uso de materiales sostenibles frente al diseño actual en el desempeño energético del edificio de la Defensoría del Pueblo en Galápagos? Para ello, se compara el escenario actual del proyecto con escenarios optimizados mediante modelado digital, simulación energética y análisis cuantitativo de indicadores.

El planteamiento metodológico toma como referencia la lógica desarrollada por \textcite{torres-barriusoSustainableBuildingStandards2025} para la definición, simulación e implementación de estándares sostenibles en Galápagos. En el artículo de referencia, el proceso parte de la caracterización de casos representativos, establece una línea base, simula estrategias sostenibles y traduce los resultados en criterios replicables. En esta tesis, dicho enfoque se aplica al proyecto de la Defensoría del Pueblo como caso de estudio en etapa de diseño, con el fin de evaluar su desempeño antes de su construcción.

La investigación es no experimental, cuantitativa, prospectiva y comparativa. Es no experimental porque no modifica físicamente el edificio; cuantitativa porque analiza indicadores medibles de desempeño energético y confort; prospectiva porque evalúa escenarios de mejora antes de su implementación; y comparativa porque contrasta el escenario base con escenarios optimizados.

El enfoque se basa en la evaluación de escenarios mediante herramientas BIM y simulación energética. La investigación no busca medir un edificio en operación, sino estimar el comportamiento energético de un proyecto arquitectónico a partir de información técnica disponible, datos climáticos y modelos de simulación.
\end{textoDosTerciosGrafico}

\clearpage
\vspace*{1.15cm}
\section{Caso de estudio y unidad de análisis}

\begin{textoDosTercios}
La unidad de análisis corresponde al proyecto arquitectónico de la Defensoría del Pueblo en la isla San Cristóbal, Galápagos. El caso se seleccionó como estudio analítico de una edificación pública ubicada en un contexto insular de alta fragilidad ambiental, donde las decisiones de diseño deben responder a condiciones climáticas, logísticas y ecológicas particulares.

El modelo base corresponde al diseño actual del edificio, incluyendo su geometría, orientación, sistema constructivo, materiales, envolvente, aperturas, condiciones de ventilación, ocupación y sistemas previstos.

La unidad de análisis se delimita al desempeño ambiental y energético del proyecto, por lo que no se evalúa la operación real de un edificio construido. Esta decisión permite trabajar con información de diseño y comparar alternativas bajo condiciones equivalentes de simulación.
\end{textoDosTercios}

\clearpage
\vspace*{1.15cm}
\section{Fases metodológicas}

\tercioTextoTikz{%
La metodología se organiza en seis fases relacionadas. Cada una transforma los resultados de la anterior en insumos para la siguiente, desde la recopilación de información y la configuración de la matriz hasta la formulación de recomendaciones.

La lectura secuencial expresa que el proceso mantiene trazabilidad entre los datos de entrada, la construcción del modelo, la simulación de escenarios y la interpretación de resultados. También permite regresar a fases previas cuando la validación exige ajustar parámetros o supuestos.
}{\graficoFasesMetodologicas}{Secuencia y productos de las seis fases metodológicas}{Diagrama original de la investigación.}

\clearpage
\vspace*{1.15cm}
\section{Fuentes de información y datos de entrada}

\begin{textoDosTercios}
La investigación utilizó información técnica del proyecto y del contexto climático. Los datos de entrada se agrupan en tres categorías: documentación del edificio, datos climáticos y parámetros de operación.

\subsection{Documentación del edificio}

La documentación utilizada comprende el modelo arquitectónico de referencia en Revit y los archivos del modelo energético. Las propiedades térmicas proceden de las construcciones y materiales registrados en el IDF. La correspondencia con las especificaciones definitivas y el juego coordinado de planos debe comprobarse antes de la ejecución; no se atribuye una revisión integral a documentos cuya validación no está registrada.
\end{textoDosTercios}

\clearpage
\vspace*{1.15cm}
\subsection{Datos climáticos}

\noindent
\begin{minipage}[t]{\textwidth}
\vspace{0pt}
\rmfamily\fontsize{12}{18}\selectfont
Se utilizó como insumo climático principal el archivo climático horario EPW de San Cristóbal, correspondiente a la estación San Cristóbal International Airport y construido como año meteorológico típico a partir del periodo 2011--2025. Este archivo contiene 8760 registros horarios y permite configurar las condiciones exteriores del modelo en DesignBuilder/EnergyPlus.

La lectura preliminar del EPW muestra una temperatura media anual de $24,28\,^{\circ}\mathrm{C}$, con valores horarios entre $15,70\,^{\circ}\mathrm{C}$ y $32,90\,^{\circ}\mathrm{C}$. La humedad relativa media es de $81,37\,$\%, la radiación global horizontal anual alcanza $2.150,50\,\mathrm{kWh}/\mathrm{m}^{2}$, equivalente a $5,89\,\mathrm{kWh}/\mathrm{m}^{2}$ por día, y la velocidad media del viento es de $4,90\,\mathrm{m/s}$. Estos valores confirman que el caso de estudio debe evaluarse como una edificación ubicada en un clima cálido-húmedo, con alta incidencia solar y potencial de ventilación natural.
\end{minipage}\par\vspace{0.5cm}
\noindent\begin{minipage}[t]{\textwidth}
\vspace{0pt}
  \captionof{table}{Variables climáticas consideradas para el modelo}
  {\fontsize{8.4}{10.2}\selectfont
  \begin{tabular}{p{0.29\linewidth}p{0.17\linewidth}p{0.44\linewidth}}
    \toprule
    Variable & Unidad & Uso metodológico \\
    \midrule
    Temperatura de bulbo seco & $^{\circ}\mathrm{C}$ & Condición exterior para cargas térmicas y confort. \\
    Punto de rocío & $^{\circ}\mathrm{C}$ & Estimación de humedad relativa y condición higrotérmica. \\
    Radiación directa y difusa & $\mathrm{kWh}/\mathrm{m}^{2}$ & Evaluación de ganancias solares y control de sombreamiento. \\
    Velocidad y dirección del viento & $\mathrm{m/s}$ y grados & Análisis del potencial de ventilación natural. \\
    Presión atmosférica & Pa & Condición atmosférica de entrada para simulación. \\
    Altura y acimut solar & grados & Definición de orientación, exposición y protecciones solares. \\
    \bottomrule
  \end{tabular}
  }
  \fuente{Basado en el archivo climático EPW de San Cristóbal.}
\end{minipage}

\clearpage
\vspace*{1.15cm}
\subsection{Condiciones de uso y ocupación}

\noindent
\begin{minipage}[t]{\textwidth}
\vspace{0pt}
\rmfamily\fontsize{12}{18}\selectfont
Se configuraron ocupación, horarios, cargas internas, iluminación, ventilación y climatización en el modelo. Son supuestos de diseño registrados en los archivos de simulación y se conservaron entre escenarios salvo las variables de intervención declaradas. No se dispone de consumos medidos para calibrar el modelo; la carga continua del datacenter requiere contraste con su inventario real.
\end{minipage}\par\vspace{0.5cm}
\noindent\begin{minipage}[t]{\textwidth}
\vspace{0pt}
  \captionof{table}{Parámetros operativos considerados en el modelo}
  {\fontsize{8.4}{10.2}\selectfont
  \begin{tabular}{p{0.24\linewidth}p{0.34\linewidth}p{0.32\linewidth}}
    \toprule
    Parámetro & Información requerida & Aplicación en la simulación \\
    \midrule
    Ocupación & Número de usuarios y horarios de permanencia. & Definir ganancias internas y patrones de uso. \\
    Iluminación artificial & Potencia instalada y horario de encendido. & Estimar cargas internas y consumo eléctrico. \\
    Equipos & Potencia de equipos de oficina y operación diaria. & Incorporar ganancias térmicas internas. \\
    Ventilación & Tipo de ventilación y disponibilidad de aperturas. & Evaluar renovación de aire y enfriamiento pasivo. \\
    Climatización & Presencia o ausencia de sistemas activos. & Comparar demanda energética y confort térmico. \\
    \bottomrule
  \end{tabular}
  }
  \fuente{Construcción metodológica a partir de los datos del estudio.}
\end{minipage}

\clearpage
\vspace*{1.15cm}
\section{Herramientas de simulación y análisis}

\begin{textoDosTercios}
El trabajo utiliza Revit como referencia geométrica y DesignBuilder para construir y simular el modelo energético mediante EnergyPlus. Las exportaciones se procesan en tablas y gráficos reproducibles, conservando sus unidades, resolución temporal y procedencia. No se atribuye validación cruzada a otros programas cuando no existe una corrida documentada en ellos.

La matriz registró los indicadores de materiales sostenibles y las evidencias necesarias para evaluarlos. Las simulaciones compararon propiedades térmicas de la envolvente; no se calculó una huella de carbono de productos sin cantidades y declaraciones ambientales verificables.

\paragraph{Construcción del caso base.}

El modelado del edificio base consiste en construir una representación digital del diseño actual del proyecto. Este modelo definió la geometría, orientación, envolvente, aperturas, sombras, materiales y condiciones operativas del edificio.

En esta fase se configuraron los principales parámetros térmicos:
\end{textoDosTercios}

\begin{textoDosTercios}
\begin{itemize}
  \item Coeficiente de transmisión térmica de muros, techos, pisos y ventanas.
  \item Espesores, densidades, conductividad y capacidad térmica de los materiales.
  \item Factores de ganancia solar de ventanas y vidrios.
  \item Hermeticidad y tasa de infiltración de aire.
  \item Temperaturas internas de referencia para calefacción o refrigeración.
  \item Humedad relativa interior y exterior.
  \item Ventilación natural y mecánica, caudal y frecuencia.
  \item Radiación solar incidente y sombreamiento existente.
  \item Sistemas de ventilación, extracción y renovación de aire.
  \item Iluminación artificial y aporte térmico asociado.
  \item Ocupación y horarios de uso del edificio.
\end{itemize}
\end{textoDosTercios}

\clearpage
\vspace*{1.15cm}
\section{Escenarios de simulación (CB, C1, C2, C3)}

\begin{textoDosTercios}
La evaluación se organizó en cuatro escenarios y cinco corridas comparativas: CB, C1, C2, C3 revisión 01 y C3 revisión 02. Las dos revisiones conservan el identificador C3 y documentan la evolución del paquete integrado. El objetivo es comparar el diseño actual con alternativas de mejora y determinar qué tipo de estrategia produce mayor incidencia sobre el desempeño energético y el confort térmico.

La comparación por escenarios permite aislar el efecto de cada grupo de decisiones y reconocer si la mejora proviene principalmente del diseño pasivo, de la materialidad o de la combinación de ambas estrategias.
\end{textoDosTercios}

\begin{textoDosTercios}
\begin{enumerate}
  \item Escenario base: diseño y materiales actuales del proyecto.
  \item C1: ventilación natural controlada en seis oficinas y aleros localizados; se conservan orientación, geometría general e inercia del caso base.
  \item C2: reducción de la absortancia solar del acabado de cubierta de 0,40 a 0,20. Se evalúa su efecto térmico; no se acredita bajo carbono ni contenido reciclado.
  \item C3: combinación inicial de C1 y C2 (revisión 01), ampliada con aislamiento, vidrio de control solar e iluminación eficiente (revisión 02). El alcance ambiental de los materiales permanece por verificar.
\end{enumerate}
\end{textoDosTercios}

\begin{textoDosTercios}
Cada escenario se modela de forma independiente y bajo criterios comparables. El alcance previsto para C2 y C3 incluye materiales sostenibles; las corridas iniciales solo ensayan un acabado reflectante, por lo que su verificación ambiental permanece pendiente (sección~\ref{sec:control-metodologico}). La propuesta busca mejorar el desempeño entre alternativas técnicamente viables; las corridas disponibles no demuestran un óptimo global.

\paragraph{Ejecución de simulaciones energéticas.}

Las simulaciones evaluaron intercambios térmicos y demanda de enfriamiento bajo el mismo clima y calendario. Se procesaron temperaturas operativas, incluida una comparación de horas ocupadas entre CB y C3 revisión 02. La luz natural disponible corresponde a un cálculo estático del caso base; no demuestra el comportamiento lumínico de todas las alternativas.

El escenario base estableció la línea de comparación. Los escenarios optimizados permitieron medir el impacto de estrategias pasivas, materiales sostenibles y combinación de ambas. Esta comparación permitió pasar de una valoración cualitativa de las estrategias a una evaluación cuantitativa de su incidencia.
\end{textoDosTercios}

\clearpage
\vspace*{1.15cm}
\section{Matriz de indicadores mínimos de sostenibilidad}
\label{sec:matriz-sostenibilidad}
\subsection{Construcción, selección y alcance de la matriz}
La matriz operacionaliza el primer objetivo específico: reúne criterios de energía, envolvente, confort, agua, materiales y operación pertinentes para un edificio público en Galápagos. La selección responde al clima cálido-húmedo, la permanencia de usuarios, la disponibilidad de agua y las exigencias de abastecimiento y mantenimiento insular. Se incluyen indicadores medibles mediante simulación y otros que requieren planos, fichas, cómputos o verificaciones posteriores. La falta de datos no elimina un indicador relevante.

La integración se realiza en cuatro pasos: identificar el criterio y su versión de origen; comprobar su pertinencia para oficinas y clima local; definir unidad, cálculo, referencia y evidencia; y aplicar la misma ficha al CB y a cada alternativa. Cuando dos marcos tratan el mismo fenómeno, se conserva un indicador y se registran las correspondencias, sin sumar puntos duplicados. La matriz es una adaptación académica mínima para este caso; su transferencia a otros edificios requiere revisar uso, clima y normativa.

Se distinguen tres referencias: el \textbf{CB de diseño}, que permite medir cambios entre escenarios; el \textbf{requisito externo}, con el alcance de su propio sistema; y el \textbf{criterio de investigación}, adoptado explícitamente cuando no se aplica un umbral externo. Una reducción frente al CB no demuestra, por sí sola, el cumplimiento de un estándar. La matriz no reemplaza la normativa local ni constituye un sistema de certificación combinado.

\clearpage
\subsection{Documentos de referencia y control de versiones}
Se adopta Minergie LATAM 2025.1 y su guía para la variante oficinas. Su certificación exige todos los requisitos obligatorios y al menos un tercio de los electivos aplicables; los indicadores seleccionados aquí no reproducen el reglamento completo \parencite{minergieLatamReglamento2025,minergieLatamGuia2025}.

Para EDGE se utilizan las guías V3 de certificación y materiales consultadas el 12 de septiembre de 2026. Se distinguen las categorías de energía operacional, agua y materiales. La guía material utiliza carbono incorporado; no se intercambian resultados de carbono y energía incorporada \parencite{ifcEdgeCertificacion2026,ifcEdgeMateriales2025}. La referencia EDGE debe construirse en su herramienta y no se presume idéntica al CB de DesignBuilder.

Para WELL se fija la edición v2 Q4 2020, características T01 y X06; se utiliza como referencia de diseño, sin afirmar adhesión a todas las actualizaciones posteriores. Para iluminación anual se fija LEED v4.1 BD+C, guía de abril de 2019, crédito EQ Daylight, opción 1. Un crédito opcional no se presenta como requisito universal de certificación \parencite{iwbiWellV22020,usgbcLeedV41Guide2019}. CEELA aporta sus 15 principios EECA como marco de diseño y operación; no se le asigna una puntuación de certificación \parencite{ceelaPrincipiosEECA2024}.

\subsection{Reglas de evaluación}
Cada fila registra evidencia, resultado y acción. Los estados son: \textbf{C}, criterio acreditado con evidencia compatible; \textbf{NC}, criterio evaluable no alcanzado; \textbf{EP}, evidencia parcial que no permite decidir; \textbf{NE}, no evaluable con la documentación disponible; \textbf{LB}, valor de línea base sin juicio de cumplimiento; y \textbf{NA}, no aplicable con justificación. Un requisito con varias condiciones solo recibe C cuando todas están acreditadas. La ausencia de datos se conserva como EP o NE y no se convierte en incumplimiento ni en cero consumo.

No se calcula un porcentaje global de sostenibilidad: los indicadores tienen alcances distintos y no se han validado ponderaciones. La cobertura documental y el cumplimiento se informan por separado. Los criterios condicionales requieren comprobar antes su aplicabilidad. Para decidir C o NC frente a una referencia térmica o energética se deben cerrar las incidencias relevantes de cargas, geometría y sistemas del modelo.

\clearpage
\subsection{Energía y envolvente}
\begin{table}[H]\centering\caption{Indicadores mínimos: energía y envolvente}
{\fontsize{9.5}{12}\selectfont\singlespacing\renewcommand{\arraystretch}{1.25}\setlength{\tabcolsep}{4pt}\begin{tabular}{>{\raggedright\arraybackslash}p{2.9cm}>{\raggedright\arraybackslash}p{5.2cm}>{\raggedright\arraybackslash}p{5.9cm}}\toprule Indicador & Criterio de referencia & Cálculo, alcance y fuente \\ \midrule
E01 · Demanda de enfriamiento & Reducir respecto al CB; criterio comparativo propio, sin umbral de certificación. & Integrar potencia horaria por zona. Ahorro = 100 × (CB $-$ caso) / CB. Mismas zonas, clima, calendario y cargas; no sumar enfriamiento sensible y total.\par\textit{CEELA, EECA 2 y 13; indicador propio} \\[4pt]
E02 · Eficiencia energética operacional & Ahorro $\geq$ 20\% en la categoría energética EDGE. & Comparar energía operacional del diseño con el caso de referencia calculado en EDGE. El CB de la tesis no sustituye automáticamente la referencia EDGE.\par\textit{EDGE V3, Part 1, p. 7; referencia externa} \\[4pt]
E03 · Transmitancia de cerramientos & U $\leq$ límite A2 aplicable. Ejemplo sin atenuación, zonas 0–2: cubierta ligera 0,40; muro sólido 0,75. & U del sistema completo y límite por elemento, clima, tipología constructiva y compacidad. Confirmar zona climática y compacidad; Ug y proporción acristalada se verifican por fachada.\par\textit{Minergie LATAM 2025.1, A2, tabla 2, pp. 15–16} \\[4pt]
E04 · Protección solar eficaz & En zonas 0–3: $\geq$ 90\% con SHGC modificado < 0,20 durante riesgo de sobrecalentamiento. & Área transparente eficazmente protegida / área transparente total. Revisar orientación, sombras y vidrio conjuntamente; longitud del alero no acredita cumplimiento.\par\textit{Minergie LATAM, guía 2025.1, A5, pp. 35–38; CEELA, EECA 2} \\[4pt]
\bottomrule\end{tabular}}\fuente{Elaboración propia con las versiones documentales indicadas en esta sección. Los criterios propios se identifican expresamente.}\end{table}\clearpage
\subsection{Eficiencia y ambiente interior}
\begin{table}[H]\centering\caption{Indicadores mínimos: eficiencia y ambiente interior}
{\fontsize{9.5}{12}\selectfont\singlespacing\renewcommand{\arraystretch}{1.25}\setlength{\tabcolsep}{4pt}\begin{tabular}{>{\raggedright\arraybackslash}p{2.9cm}>{\raggedright\arraybackslash}p{5.2cm}>{\raggedright\arraybackslash}p{5.9cm}}\toprule Indicador & Criterio de referencia & Cálculo, alcance y fuente \\ \midrule
E05 · Potencia y control de iluminación & $\leq$ 8 W/m\textsuperscript{2}, LED y control sobre $\geq$ 80\% de potencia o todos los recintos regularmente ocupados. & Potencia instalada / área por recinto; verificar LED y alcance de automatización. Requisito para oficinas y espacios exteriores; conservar iluminancia de tarea y comprobar LED con ficha técnica.\par\textit{Minergie LATAM 2025.1, T3, p. 27; CEELA, EECA 10} \\[4pt]
C01 · Aceptabilidad térmica ocupada & Ruta mecánica seleccionada: PMV entre $-$0,5 y +0,5 en $\geq$ 90\% de horas y $\geq$ 90\% de espacios regularmente ocupados. & Evaluar PMV por recinto y hora; en modo natural aplicar modelo adaptativo admisible. En modo mixto separar periodos mecánicos y naturales; documentar met, clo, velocidad, humedad y radiación.\par\textit{WELL v2 Q4 2020, T01.1, opción 1, p. 161} \\[4pt]
C02 · Suministro de aire exterior & En los recintos indicados por T4: caudal $\geq$ 1,30 veces referencia local o ASHRAE 62.1 si no existe. & Separar aire exterior de infiltración; verificar caudal, filtración y recuperación. Aplicar variante oficinas y condiciones de exención; ventilación natural no prueba calidad de aire.\par\textit{Minergie LATAM 2025.1, T4, pp. 28–29; CEELA, EECA 6} \\[4pt]
C03 · Autonomía lumínica y exposición solar & sDA $\geq$ 40\% como nivel de entrada del crédito; documentar control de deslumbramiento donde ASE supere 10\%. & Calcular sDA300/50\% y ASE1000,250 con simulación anual. Crédito opcional usado como referencia; no es requisito mínimo para certificar todo el edificio.\par\textit{LEED v4.1 BD+C, guía abril 2019, EQ Daylight, opción 1} \\[4pt]
\bottomrule\end{tabular}}\fuente{Elaboración propia con las versiones documentales indicadas en esta sección. Los criterios propios se identifican expresamente.}\end{table}\clearpage
\subsection{Agua, materiales y generación}
\begin{table}[H]\centering\caption{Indicadores mínimos: agua, materiales y generación}
{\fontsize{9.5}{12}\selectfont\singlespacing\renewcommand{\arraystretch}{1.25}\setlength{\tabcolsep}{4pt}\begin{tabular}{>{\raggedright\arraybackslash}p{2.9cm}>{\raggedright\arraybackslash}p{5.2cm}>{\raggedright\arraybackslash}p{5.9cm}}\toprule Indicador & Criterio de referencia & Cálculo, alcance y fuente \\ \midrule
A01 · Uso eficiente del agua & Ahorro $\geq$ 20\% en la categoría agua EDGE. & Volumen de agua según usos, caudales, frecuencias y ocupación; contraste EDGE. Mantener igual población y frecuencia de uso; documentar todos los usos aplicables.\par\textit{EDGE V3, Part 1, p. 7; CEELA, EECA 12} \\[4pt]
M01 · Carbono incorporado & Meta de reducción $\geq$ 20\% para categoría materiales EDGE; en V3, cálculo de carbono incorporado. & Sumar cantidades × factores ambientales compatibles; contraste material EDGE. Fijar misma unidad funcional y alcance de módulos; añadir transporte insular por separado.\par\textit{EDGE V3, Part 5, introducción y anexo 1; CEELA, EECA 3} \\[4pt]
M02 · Emisiones de productos interiores & Meta documental propia: 100\% de productos de la intervención con verificación de emisiones; aplicar límites del ensayo seleccionado. & Inventario de acabados, adhesivos y pinturas con ficha de emisiones aplicable. La cobertura documental no equivale a obtener el crédito WELL X06.\par\textit{WELL v2 Q4 2020, X06; CEELA, EECA 5} \\[4pt]
E06 · Autoproducción & Ruta fotovoltaica: $\geq$ 0,010 kWp/m\textsuperscript{2} SRE; el requisito se limita a 30 kWp. & Dimensionar respecto a superficie de referencia energética y estimar producción. Ruta seleccionada para oficinas; justificar superficie, sombras y posibles restricciones locales.\par\textit{Minergie LATAM 2025.1, T2, p. 26; CEELA, EECA 14} \\[4pt]
\bottomrule\end{tabular}}\fuente{Elaboración propia con las versiones documentales indicadas en esta sección. Los criterios propios se identifican expresamente.}\end{table}\clearpage
\subsection{Sistemas, operación y adaptación insular}
\begin{table}[H]\centering\caption{Indicadores mínimos: sistemas, operación y adaptación insular}
{\fontsize{9.5}{12}\selectfont\singlespacing\renewcommand{\arraystretch}{1.25}\setlength{\tabcolsep}{4pt}\begin{tabular}{>{\raggedright\arraybackslash}p{2.9cm}>{\raggedright\arraybackslash}p{5.2cm}>{\raggedright\arraybackslash}p{5.9cm}}\toprule Indicador & Criterio de referencia & Cálculo, alcance y fuente \\ \midrule
E07 · Refrigeración eficiente & SEER $\geq$ 6 kW/kW (método europeo) y sobredimensionamiento $\leq$ 15\% en oficinas; revisar refrigerante y recuperación. & Verificar ficha de equipo, refrigerante y potencia instalada frente a carga de diseño. No convertir directamente COP del modelo en SEER; aplicar las condiciones climáticas de T5.\par\textit{Minergie LATAM 2025.1, T5, p. 30; CEELA, EECA 13} \\[4pt]
G01 · Monitoreo y control & Plan completo de energía, agua y ambiente interior para la fase de operación. & Comprobar plan con variables, medidores, responsables, frecuencia y mantenimiento. Evaluación de diseño; los datos simulados no sustituyen medición posterior.\par\textit{CEELA, EECA 15; Minergie O1–O3; criterio documental propio} \\[4pt]
G02 · Viabilidad de materiales y soluciones & Todas las soluciones propuestas con ficha de viabilidad y fuentes; menor impacto debe demostrarse. & Comparar procedencia, transporte, vida útil, humedad/salinidad, mantenimiento y costo. Criterio propio de selección; no supone que cualquier material local o reciclado sea automáticamente sostenible.\par\textit{CEELA, EECA 1 y 3; adaptación propia a Galápagos} \\[4pt]
P01 · Apertura útil y enfriamiento natural & Si se emplea enfriamiento nocturno natural: aberturas $\geq$ 4\% del piso en oficinas, protegidas de lluvia y robos. & Abertura neta disponible / área de piso del recinto; comprobar recorrido y horario. Criterio condicionado a esa estrategia; fracción practicable de ventana y proporción respecto al piso son diferentes.\par\textit{Minergie LATAM 2025.1, A6, pp. 18–19; CEELA, EECA 6 y 8} \\[4pt]
\bottomrule\end{tabular}}\fuente{Elaboración propia con las versiones documentales indicadas en esta sección. Los criterios propios se identifican expresamente.}\end{table}\clearpage
\subsection{Aplicación y validación comparativa}
La matriz se aplica primero al caso base (sección~\ref{sec:matriz-caso-base}) y después a las alternativas (sección~\ref{sec:matriz-escenarios}). Para cada corrida se conservan clima, calendario, ocupación, cargas de referencia y alcance de zonas. Las variables de intervención se registran antes de interpretar resultados. Si una revisión modifica un supuesto del CB, deben recalcularse las alternativas que se comparen con esa versión.

Para una magnitud cuyo descenso es favorable, la mejora porcentual se define como $100(X_{CB}-X_i)/X_{CB}$. En confort se reporta también la diferencia en puntos porcentuales de horas aceptables: $H_i-H_{CB}$. Si el denominador es cero o faltan datos, no se informa un porcentaje. Los porcentajes mensuales no se promedian para obtener el anual.

La evaluación térmica se realiza por recinto regularmente ocupado y horario efectivo. En espacios de modo mixto se aplica el método correspondiente al modo activo en cada periodo. Se documentan vestimenta, actividad, humedad, temperatura radiante y velocidad del aire. Las medias del edificio y el recuento de días sobre 28 °C se conservan solo como diagnóstico exploratorio. No se equipara una reducción de enfriamiento a mejora del confort.

La comparación de materiales exige igual función, vida útil y límites del cálculo ambiental. Una mejora de reflectancia o transmitancia puede reducir demanda térmica, pero no demuestra menor carbono incorporado. La validación de diseño comprende consistencia del modelo, trazabilidad de datos y sensibilidad de supuestos; sin mediciones de operación no se declara calibración empírica. La revisión documental de la matriz queda abierta a contraste técnico externo antes de proponer su uso generalizado.

\clearpage
\subsection{Trazabilidad entre objetivos, método y productos}
\label{sec:trazabilidad-objetivos}
Los cuatro objetivos se articulan en una secuencia: configurar la matriz, aplicarla al diseño inicial, formular una propuesta que responda al diagnóstico y validar sus efectos. Simular casos no sustituye la aplicación de los indicadores ni resuelve automáticamente el desempeño ambiental de los materiales.

\begin{table}[H]\centering
\caption{Correspondencia entre objetivos específicos y productos verificables}
{\small\singlespacing\renewcommand{\arraystretch}{1.35}\setlength{\tabcolsep}{4pt}
\begin{tabular}{>{\raggedright\arraybackslash}p{2.6cm}>{\raggedright\arraybackslash}p{5.6cm}>{\raggedright\arraybackslash}p{5.8cm}}\toprule
Objetivo & Método & Producto y criterio de cierre \\ \midrule
OE1. Configurar la matriz & Selección, adaptación y documentación de indicadores de los cinco marcos. & Capítulo 3: matriz con fuente, versión, unidad, cálculo, referencia y regla de decisión. \\
OE2. Evaluar el CB & Aplicación de cada indicador al modelo y documentación del diseño inicial. & Capítulo 4: resultados, brechas y no evaluables por indicador, sin confundir línea base y certificación. \\
OE3. Desarrollar la propuesta & Medidas pasivas y materiales seleccionados a partir de las brechas, con viabilidad documentada. & Capítulo 5: escenarios trazables y especificaciones técnicas y ambientales justificadas. \\
OE4. Validar la incidencia & Comparación energética y de confort en condiciones equivalentes; revisión de límites e incertidumbres. & Capítulo 5: diferencias de energía y de horas aceptables por zona; síntesis de alcance en conclusiones. \\
\bottomrule\end{tabular}}
\fuente{Elaboración propia a partir de los objetivos específicos de la investigación.}
\end{table}

\clearpage
\section{Procedimiento y validación del modelo}

\begin{textoDosTercios}
El procedimiento operativo se desarrolló en las siguientes etapas:

Estas etapas ordenan el proceso de trabajo para evitar que los resultados dependan de decisiones aisladas o de cambios no controlados entre modelos.
\end{textoDosTercios}

\begin{textoDosTercios}
\begin{enumerate}
  \item Configuración de la matriz con fuentes, versiones, criterios y reglas de evaluación.
  \item Revisión del modelo arquitectónico disponible y delimitación de información documental pendiente.
  \item Recopilación de especificaciones técnicas de materiales y sistemas actuales.
  \item Recolección de datos climáticos de San Cristóbal, Galápagos.
  \item Registro de supuestos de uso y de las condiciones que requieren confirmación con responsables del proyecto.
  \item Creación del modelo digital del edificio en software especializado.
  \item Definición de parámetros térmicos, propiedades de materiales y sistemas HVAC.
  \item Configuración de condiciones climáticas locales y horarios de ocupación.
  \item Aplicación de la matriz al caso base e identificación de brechas prioritarias.
  \item Diseño de escenarios alternativos con estrategias pasivas y materiales sostenibles que respondan al diagnóstico.
  \item Ejecución de simulaciones energéticas para cada escenario.
  \item Registro sistemático de resultados para análisis comparativo.
  \item Procesamiento de resultados mediante tablas, gráficos y matrices.
  \item Identificación de estrategias y materiales con mayor impacto.
  \item Identificación de requisitos de viabilidad técnica, suministro y costo para una evaluación posterior.
  \item Formulación de recomendaciones de optimización según impacto y factibilidad.
\end{enumerate}
\end{textoDosTercios}

\begin{textoDosTercios}
La evaluación económica y el análisis de ciclo de vida quedan fuera del alcance ejecutado. No se calcularon costos de inversión, retorno económico ni ahorro de carbono incorporado.

\paragraph{Producto final de la metodología.}

El producto final de la metodología es una propuesta de optimización del diseño arquitectónico mediante estrategias pasivas y materiales sostenibles. Esta propuesta se apoya en evidencia cuantitativa derivada de simulaciones, comparación de escenarios e identificación de prioridades de implementación.
\end{textoDosTercios}
